\chapter{Mantle flow and its kernels}
\label{ch:flow}

This chapter explains the part of the model that everything else rests on: how a density anomaly in the mantle produces a stress on the base of the lithosphere and a change in the gravitational potential.
It is the part you wrote and understand best, so the aim here is to make explicit the choices that are easy to take for granted.

\section{The equations}

The mantle is treated as a viscous fluid at infinite Prandtl number: inertia is negligible, so at every instant the forces balance.
We solve for the instantaneous flow driven by a density anomaly $\delta\rho_m$, with the viscosity $\eta(r)$ a function of radius only.
The governing equations, linearised about a hydrostatic, spherically symmetric reference state $\rho_0(r)$, $g_0(r)$, are
\begin{align}
    \nabla\cdot\bm{\sigma} - \delta\rho\, g_0(r)\,\hat{\bm r} - \rho_0\nabla\phi_1 &= 0
    && \text{(momentum: force up = force down)}, \label{eq:f-mom}\\
    \nabla\cdot(\rho_0\bm u) &= 0
    && \text{(mass, anelastic)}, \label{eq:f-mass}\\
    \nabla^2\phi_1 &= 4\pi G\,\delta\rho
    && \text{(Poisson, physics sign)}, \label{eq:f-poisson}
\end{align}
with $\bm\sigma = -p\bm I + \eta(\nabla\bm u + \nabla\bm u^{\mathsf T} - \tfrac23(\nabla\cdot\bm u)\bm I)$, stress positive in tension, and $\delta\rho$ the density anomaly.
Compressibility enters through the anelastic mass equation, $\nabla\cdot(\rho_0\bm u)=0$, and the $-\tfrac23(\nabla\cdot\bm u)$ term in the stress; $\bm u$ is a velocity, so there is no advected density term in an instantaneous problem. (Check the exact form flowkernels solves before quoting this equation in the thesis.)

Reading Equation~\eqref{eq:f-mom} as a force balance: the divergence of stress (the net viscous and pressure force on a parcel) holds up the parcel's extra weight $\delta\rho\,g_0$ and its weight in the perturbed gravity field $\rho_0\nabla\phi_1$.
The last term is self-gravitation: the potential of the anomaly and of the boundaries it deflects acts back on the reference density.

\begin{decision}{Compressible and self-gravitating}
\textbf{Choice.} Keep the anelastic compressibility term and the self-gravitation term $\rho_0\nabla\phi_1$, following \citet{corrieu_mantle_1995}.

\textbf{Alternatives.} The incompressible, non-self-gravitating Stokes problem of \citet{hager_lower_1985} and many later studies.

\textbf{Why.} Both effects change the geoid kernel appreciably at low degree, which is exactly where the geoid has its power \citep{corrieu_mantle_1995,forte_viscous_1991}; self-gravitation matters most at degrees 2--3.
(Quantify this yourself before quoting a number: run flowkernels with each term switched off.)
Since the cost of including them in a sparse solver is negligible, there is no reason to leave them out.

\textbf{Cost.} None numerically; the equations are longer, and comparison with incompressible studies needs care.
\end{decision}

\section{Boundary conditions}

\paragraph{Free slip at the core--mantle boundary.}
The outer core is an inviscid fluid, so it exerts no shear stress on the mantle: $\sigma_{r\theta}=\sigma_{r\varphi}=0$ and $u_r = 0$ at $r_{\mathrm{cmb}}$.
The normal stress there is not zero; it is what deflects the CMB, and the deflected CMB is a mass anomaly that contributes to the geoid.

\paragraph{Free slip at the top of the flow, $r_{\mathrm{top}} = R - \qty{105}{km}$.}
The top of the flow domain is the base of the lithosphere.
We impose $u_r=0$ and zero shear there, and record the normal stress $\sigma_{rr}(r_{\mathrm{top}})$: that is the stress the mantle exerts on the base of the lithosphere, which the lithosphere then has to balance.
This is the classical way of computing dynamic topography \citep{hager_lower_1985,richards_geoid_1984}: solve with an undeformed boundary, then convert the normal stress into a deflection.
It is exact to first order because the deflection is small compared with the wavelengths of interest.

\begin{decision}{The top of the flow at \qty{105}{km}}
\textbf{Choice.} The flow domain stops at \qty{105}{km} depth, and everything above is ``lithosphere'', which is carried by the flow and loaded by it but does not flow.

\textbf{Alternatives.} Flow to the surface (as in \citet{hager_lower_1985}), with lateral viscosity variations representing plates; or a deeper base (\qtyrange{150}{200}{km}).

\textbf{Why.}
\begin{itemize}[nosep]
\item Our observables are corrected for, or later model explicitly, the lithosphere, so the mantle model must stop where the lithosphere starts.
\item The area-weighted mean depth of the WINTERC-G lithosphere--asthenosphere boundary is \qty{102}{km} and its median \qty{87}{km}, so \qty{105}{km} is a representative depth.
\item The flowkernels grid has a node close to \qty{105}{km}, and the core model and the oceanic residual depths it is compared with are both defined relative to a lithosphere of roughly this thickness.
\end{itemize}
\textbf{Caveat.} About \qty{28}{\percent} of the Earth's area has a WINTERC-G LAB deeper than \qty{105}{km}, mostly beneath cratons.
There, the lithospheric mantle below \qty{105}{km} is treated as mantle and its density as a mantle anomaly.
This is consistent (it is within the flow domain and moves with it) but means the deepest lithosphere is not part of the lithospheric unknowns of Chapter~\ref{ch:lithosphere}.
In the lid-base table of Chapter~\ref{ch:lithosphere}, moving the base to \qtyrange{150}{260}{km} changes the reference topography substantially, which is why the choice is stated explicitly.
\end{decision}

\section{Degree one and the reference frame}

At degree 1 a density anomaly displaces the centre of mass of the mantle.
Without a frame condition the problem is singular: a rigid translation of the whole Earth costs nothing.
We work in the centre-of-mass frame of the whole Earth, in which the external potential has no degree-1 part.
In flowkernels this is imposed by replacing the potential-matching row at the surface by $\Phi(b) = 0$ at degree 1.

The consequence for the observables is simple.
In this frame the degree-1 geoid is zero, $N_1=0$, but the degree-1 topography is not: the surface moves relative to the centre of mass, and that motion is a real, observable height difference between hemispheres.
The deflection $w_1$ on its own depends on the frame, but $h_1 = w_1 - N_1$, a height measured from sea level, does not; in the centre-of-mass frame $N_1=0$, so there $h_1=w_1$.
One subtlety: the zero-dipole condition is imposed in flowkernels with the deflected top's mass at $r_{\mathrm{top}}$; the core model moves that mass to $R$, where its dipole moment is larger by $(R/r_{\mathrm{top}})^3\approx1.05$, so the model's mass dipole is no longer exactly zero.

\begin{defence}
\Q{Why include degree 1 at all? Many studies drop it.}
\A Degree 1 is part of the observed topography and our model predicts it.
Dropping it throws away the hemispheric asymmetry of the data, which the inversion should be allowed to use or to fail to fit.
The degree-1 kernels of flowkernels agree with the independent finite-element solver \texttt{fem\_1d} to a relative misfit of $10^{-4}$, so the prediction is reliable.
Including it lowered the Bayes factor for topography from 39.0 to 35.2 (Chapter~\ref{ch:single}), so the degree-1 data on their own favour the null model by roughly 4 nats (roughly, because the hyperparameters are refitted): the scaling does not explain the degree-1 topography.

\Q{Why is the degree-1 geoid zero?}
\A Because we work in the centre-of-mass frame, in which the external field has no dipole by definition.
The satellite geoid is also expressed in this frame, so the observed degree-1 coefficients are zero too.
\end{defence}

\section{The viscosity profile}

The viscosity enters the kernels; the density does not depend on it, but what we infer about the density does.
The reference profile has a factor-30 increase at \qty{670}{km}, an approximation to the jump inferred from the geoid by \citet{hager_lower_1985} and later studies \citep{king_masters_1992,rudolph_2015}.
On the flowkernels grid the jump sits at the node nearest \qty{670}{km}, at \qty{666}{km}; the audit quotes an effective position of \qty{672}{km} (finding M6); check how that number was defined before quoting it.

\begin{decision}{A factor-30 jump at \qty{670}{km} as the reference}
\textbf{Choice.} Reference viscosity $\eta(r)$ uniform in the upper mantle, $\times 30$ below \qty{670}{km}.

\textbf{Alternatives.} Factors of 10 and 100; a constant viscosity; profiles with a low-viscosity asthenosphere or a viscosity hill in the mid-mantle \citep{rudolph_2015}.

\textbf{Why.} It is the most commonly used simple profile and sits in the middle of the range inferred from the geoid and glacial isostatic adjustment.
More importantly, the inversion was repeated for factors 10, 30 and 100 and for constant viscosity: the log Bayes factor against the null varied only from 30 to 35 across them (in the earlier full-model run; still to be repeated with the core model), so the data do not strongly prefer any.

\textbf{Caveat.} The recovered density does depend on the viscosity.
How strongly is what the adjoint viscosity kernels of Section~3 of the chapter quantify.
\end{decision}

\section{Why a sparse boundary-value solver}

The classical way to solve the radial equations for each degree is the propagator-matrix method: start at the CMB with unknown amplitudes, integrate the ODE system up to the surface, and match the boundary conditions there \citep{hager_lower_1985,corrieu_mantle_1995}.
Its difficulty is that the fundamental solutions grow and decay as $r^{\pm\ell}$.
Over the mantle, $r_{\mathrm{top}}/r_{\mathrm{cmb}} \approx 1.8$, so at degree 40 a growing solution increases by about $1.8^{40}\approx 10^{10}$ while a decaying one falls by about as much: the two differ by some 20 orders of magnitude, and matching boundary conditions means subtracting numbers of very different size.
The resulting linear system is exponentially ill-conditioned with degree.

flowkernels instead discretises the whole radial interval at once and solves the discrete equations, with all the boundary conditions, as one sparse linear system \citep{ascher_bvp_1995}.
Each equation only couples neighbouring nodes, so the system is banded; its conditioning grows only algebraically with the number of nodes, not exponentially with degree. The number of nodes needed for a given accuracy still grows with degree (a deep, sharp anomaly at $\ell=40$ is \qty{9}{\percent} off on the present grid; Chapter~\ref{ch:assembly}), but that is a resolution question, not a loss of precision.

\begin{defence}
\Q{Why not propagator matrices, which are standard?}
\A Because their conditioning grows exponentially with degree.
Remedies exist (renormalisation, compound matrices) but they patch a problem that a global sparse solve does not have.
We need degrees up to 40 for the topography, and the same code solves every degree with the same accuracy.

\Q{How do you know the solver is right?}
\A It agrees with \texttt{fem\_1d}, an independently written finite-element solver, to \qty{0.09}{\percent} for the topography kernels and \qty{0.1}{\percent} for the geoid kernels across degrees 1--40, and the limit tests in Chapter~\ref{ch:assembly} close to $10^{-10}$.
\end{defence}

\section{From traces to kernels}

flowkernels returns, for each degree and for a unit density sheet at each radial node, five traces: the tangential velocity at the top, the normal stress and potential at the top, and the normal stress and potential at the bottom.
From these we form the kernels.

\paragraph{The load on the plate.}
The normal stress at the top, $\sigma_{rr}(r_{\mathrm{top}})$, is non-dimensional in flowkernels; multiplying by $\eta_{\mathrm{ref}}/r_{\mathrm{top}}$ restores units of pascals.
With stress positive in tension, $\sigma_{rr}>0$ means the material on either side of a horizontal surface pulls on the other, and $\sigma_{rr}<0$ means it pushes (in a fluid at rest $\sigma_{rr}=-p$, so a positive pressure is a negative $\sigma_{rr}$).
The traction on a surface with outward normal $\hat{\bm n}$ is $\bm t = \bm\sigma\cdot\hat{\bm n}$; the base of the lithosphere faces down, $\hat{\bm n} = -\hat{\bm r}$, so the radial force per unit area on the lithosphere from the mantle is $t_r = -\sigma_{rr}$.
The mantle therefore pushes up where $\sigma_{rr}<0$ (compression, above a rising hot anomaly) and pulls down where $\sigma_{rr}>0$ (tension, above a sinking dense one).
This is the perturbation stress of the flow: the hydrostatic part that holds up the reference lithosphere has been removed.
So the upward load is
\begin{equation}
    q = -\frac{\eta_{\mathrm{ref}}}{r_{\mathrm{top}}}\,\sigma_{rr}(r_{\mathrm{top}}),
    \qquad K_q = -\frac{\eta_{\mathrm{ref}}}{r_{\mathrm{top}}}\,\sigma_{\mathrm{top}} .
\end{equation}

\paragraph{The geoid without the deflected top.}
flowkernels' own geoid at the top of the shell, $K_N^{\mathrm{fk}}$, includes the mass of the deflected top boundary, placed at $r_{\mathrm{top}}$.
In our model that mass is not at $r_{\mathrm{top}}$: the whole lithospheric column moves, and its mass change sits wherever the column's density contrasts are (at the surface in the core model).
So we remove it.
flowkernels' surface potential row reads $\Phi + c_{\mathrm{top}}\sigma_{rr} + (\text{CMB terms}) = (\text{attraction of the anomaly})$, where
\begin{equation}
    c_{\mathrm{top}} = \frac{4\pi G}{2\ell+1}\,\frac{r_{\mathrm{top}}}{g_{\mathrm{top}}}\,\rho_{\mathrm{ref}}
\end{equation}
is the self-gravity of the deflected top, written in flowkernels' own sign convention for $\Phi$ (as in the comment in \texttt{mantle.py}); its sign depends on that convention, which the limit test of Chapter~\ref{ch:geoid} confirms numerically to $10^{-10}$.
Adding $c_{\mathrm{top}}\sigma_{\mathrm{top}}$ back to $\Phi_{\mathrm{top}}$ leaves the potential of the anomaly and of the deflected CMB alone.
We then continue it upward to the surface as $(r_{\mathrm{top}}/R)^{\ell+1}$, the decay of an external potential field, and divide by surface gravity:
\begin{equation}
    K_N = \frac{\eta_{\mathrm{ref}}}{\rho_{\mathrm{ref}}\,r_{\mathrm{top}}\,g_{\mathrm{top}}}
    \left(\Phi_{\mathrm{top}} + c_{\mathrm{top}}\sigma_{\mathrm{top}}\right)
    \frac{g_{\mathrm{top}}}{g_0}\left(\frac{r_{\mathrm{top}}}{R}\right)^{\ell+1},
    \qquad K_N = 0 \text{ at } \ell=1 .
\end{equation}

\begin{decision}{Remove the top's mass from flowkernels' geoid and put it back where the column puts it}
\textbf{Choice.} Use $K_N$ (anomaly and CMB) and add the deflected column's mass separately, at the surface in the core model and at each density contrast in the full model.

\textbf{Alternative.} Use flowkernels' geoid $K_N^{\mathrm{fk}}$ directly, which puts the compensating mass at \qty{105}{km}.

\textbf{Why.} The geoid of a compensated anomaly is the small difference between the anomaly's attraction and that of its compensating mass, so it is very sensitive to where the compensating mass sits.
Using $K_N^{\mathrm{fk}}$ in the core model gave a geoid wrong by factors of 0.2 to 3 at low degree, compared with the full model.
Moving the mass to the surface removed the discrepancy (Chapter~\ref{ch:core}).
\end{decision}

\paragraph{Quadrature in radius.}
The observables are integrals over depth, Equation~\eqref{eq:p-kernels}, computed with trapezoid weights on the kernel grid.
The density, and so the integrand, can be discontinuous at the viscosity jump, so the grid is split there: the node on the jump is listed twice, once as the top of the lower layer and once, nudged up by a relative $10^{-12}$, as the bottom of the upper layer, and each layer gets its own trapezoid weights.
A density step aligned with the jump is then not smeared across it (the trapezoid rule is exact only for an integrand linear within each layer).

\begin{incode}{\texttt{src/density\_inversion/forward/mantle.py}, \texttt{MantleFlow.\_kernels}}
For each degree $\ell = 1,\dots,L$ the function calls \texttt{fk.solve\_degree\_traces(shell, l)} and builds
\begin{itemize}[nosep]
\item \texttt{K\_q} $= -(\eta_{\mathrm{ref}}/r_{\mathrm{top}})\,\sigma_{\mathrm{top}}$: the load on the plate, \si{Pa} per (\si{kg/m^3} per m of sheet);
\item \texttt{K\_h0}: flowkernels' own water-loaded free-surface topography, kept for comparison only;
\item \texttt{K\_N}: the geoid of the anomaly and the CMB, continued to the surface, zero at $\ell=1$;
\item \texttt{K\_N\_fk}: flowkernels' geoid including the top at \qty{105}{km}, kept for the limit test;
\item \texttt{K\_m\_top}: the mass per unit area of the deflected top.
\end{itemize}
It then splits the trapezoid weights at the viscosity jump (raising an error if there is none) and caches everything to an \texttt{.npz} file keyed on the viscosity, LAB depth and $L$.
The kernels are used as
\begin{lstlisting}
q_lm = sum_r weights * K_q[l-1] * drho_lm(r)
\end{lstlisting}
which is exactly the quadrature of Equation~\eqref{eq:p-kernels}.
\end{incode}
